% ============================================================
%  H. Høffding (ed.), "Udvalgte Stykker af dansk filosofisk Litteratur"
%  = Mindesmærker af Danmarks Nationallitteratur, III. Copenhagen:
%    Gyldendal 1910.
%  TRANSLATION (English) — Hans Halvorson
%
%  Høffding's editorial matter for section IV, "The Psychological School
%  (Treschow, Sibbern, Howitz)": the section heading, his introductory
%  note § 1 on Treschow, and the headings of the two Treschow excerpts.
%  Source of truth: transcription.tex. Mirrors it 1:1.
%
%  TERMINOLOGY — aligned with ../../treschow/enslige-ting/translation.tex:
%    enslige Ting -> singular things;  det Individuelle -> the individual
%    det egentligt Reale -> what is properly real
%    Uudtømmelighed -> inexhaustibility
%    Erfaring / Erfaringssynspunkter -> experience / the standpoint of experience
%  Høffding writes modern (1910) Danish; Treschow's own 1807/1810 spelling
%  is preserved only inside the quoted essay title.
% ============================================================
\documentclass[12pt,a4paper]{article}
\usepackage[utf8]{inputenc}
\usepackage[T1]{fontenc}
\usepackage{libertinus}
\usepackage{libertinust1math}
\usepackage[english]{babel}
\usepackage[protrusion=true,expansion=true]{microtype}
\usepackage{setspace}
\usepackage[margin=1.2in]{geometry}
\usepackage{fancyhdr}
\usepackage{hyperref}

\hypersetup{
  pdftitle={Selected Passages from Danish Philosophical Literature — Treschow},
  pdfauthor={Harald Høffding},
  hidelinks
}

\pagestyle{fancy}
\fancyhf{}
\rhead{Høffding (ed.), \emph{Selected Passages} (1910) --- Treschow}
\cfoot{\thepage}

\setstretch{1.3}

\title{\textbf{The Psychological School}\\[4pt]
  \large Høffding's Introduction to Treschow\\[6pt]
  \normalsize in \emph{Selected Passages from Danish Philosophical Literature}}
\author{Harald Høffding (editor)\\[4pt]
  \small Translated from the Danish by Hans Halvorson}
\date{\small Monuments of Denmark's National Literature, III.\\
  Copenhagen: Gyldendal, 1910}

\begin{document}
\maketitle
\thispagestyle{fancy}

\bigskip

% === TRANSLATION BEGINS (printed p. 99) ===

% printed p. 99
\begin{center}
  {\large IV}\\[10pt]
  {\Large THE PSYCHOLOGICAL SCHOOL (TRESCHOW,\\ SIBBERN, HOWITZ).}
\end{center}

\begin{center}\rule{0.10\linewidth}{0.4pt}\end{center}

\noindent
1.~That interest in experience, and in the problems that arise on the ground
of experience, by which Holberg and Sneedorff had been animated, found voice
again in a series of men who in the nineteenth century cultivated psychology
and sought partly to bring it into close connection with natural science,
partly to treat wider philosophical questions from the standpoint of
experience. The first in the series is \textsc{Niels Treschow} (1751--1833),
author of the first independent Danish psychology. In his philosophy he
maintained---in opposition to German philosophy's absorption into general and
abstract concepts---the individual and the concrete as what is properly real
(``Is there any concept or any idea of singular things?'' Videnskabernes
Selskabs Skrifter 1807); and it was therefore natural for him to suppose that
the differences between species had developed historically and successively.
The ``moral'' misgivings about this supposition he has refuted in a fine manner
(Elementer til Historiens Filosofi 1811).

\bigskip
\noindent
{\large A. THE INEXHAUSTIBILITY OF THE INDIVIDUAL}

% [There follows the excerpt from "Gives der noget Begreb eller nogen Idee om
%  enslige Ting?" (1810), anthology pp. 99--108. It reproduces the original's
%  printed pp. 225--230 and 235--241 (first line); pp. 231--234 and 241--254 are
%  omitted, and the omissions are NOT marked. Complete text and translation:
%  ../../treschow/enslige-ting/. The excerpt ends: "...properly, however, it is
%  right to make a distinction between this abstract representation of the
%  singular and the concrete one."]

\bigskip
\begin{center}\rule{0.22\linewidth}{0.4pt}\end{center}

% printed p. 108
\noindent
{\large B. THE DOCTRINE OF DEVELOPMENT}

% SOURCE (verified): "Elementer til Historiens Philosophie" (1811). Høffding
% gives no source at the excerpt itself; his § 1 cites the work. Verified
% against nb.no's full-text index: "Thibetanerne" and "Evolutions" occur ONLY
% in Elementer (1811) among Treschow's works, and anthology pp. 111--112 are
% word-for-word identical with Elementer, printed p. 80 (scan p. 92,
% image-verified). NB: the original is Fraktur and uses "ø"; Høffding has
% normalized the orthography. Cite the original, not the excerpt.
% Terminology: Særsyner -> singular appearances (as in the 1810 essay);
% Slægter/Arter -> genera/species; Kimer -> germs; Grundformerne -> the
% fundamental forms; Udartninger -> degenerations.

\noindent
We gladly see the singular appearances that occur in experience explained in
some natural way. Our desire for knowledge is thereby satisfied, whereas it is
rebuffed or thrust back when omnipotence or some other supernatural cause is
put forward. In the childhood of reason one does not know the distinction
between the natural and the supernatural: the two mingle in the less cultivated
person's manner of representation, and are the same. It is perhaps less a real
error than an idea not distinctly thought. But however that may be, therein
lies the whole reason why the ancients,
% printed p. 109
in all cases usual and unusual alike, turned their thoughts to the influence of
a god or of a hidden power; whereas we always believe that the nearest causes
must be found in something visible and known, and regard only the higher causes
as concealed and supersensible.

As concerns the first origin of the world, of animals and of plants, we are
therefore inclined to assume an inconceivable creation, although we grant that
in their subsequent changes and propagation nature's own forces, according to
laws once determined, are chiefly at work. If one carries the merely physical
investigations somewhat further, one does seem to become aware of a kind of
possibility of seeing how the so-called lifeless materials and bodies may have
formed themselves by an innate force; but one despairs of grasping the
production of organic bodies in the same way. From this arises the
evolution-system, adopted until the most recent times with almost universal
approval, whose ultimate ground is the many kinds of shoots or germs that
seemed able to have received existence only by an immediate creation. That these
could have been produced by a raw matter's gradually self-developing force---this
opinion seemed to most as unreasonable as it was dangerous and bordering on
explicit denial of God. How deeply, it was further thought, is human nature
thereby degraded, indeed every living being---although the ancients did not find
it degrading, but reverently worshipped the earth as their mother. Yet we see
that such is nature's course. Of omnipotence one forms the conception that it
can bring forth everything at once, the greatest along with the least. One
wonders, therefore, and rightly, on learning that everything comes to be out of
something preceding which, when we consider the whole, is always, as far as form
goes, inferior to what follows. The seed or germ, even of human beings and
animals, is in the beginning nothing but a wholly formless fluid: if we pursue
it to its origin---with our thoughts, when the eye can no longer do so---then it
loses itself at last in materials that we cannot
% printed p. 110
distinguish from the common ones. But how little are these known to us! How
glorious is that dust which we use as an emblem of lowliness and out of
ignorance despise! Is the common matter anything other than the visible God,
than a revelation of the highest being itself?\footnote{By matter one must here
understand not the phenomenon, which is an object of experience, but its
supersensible substrate, the ultimate subject of the same, its eternal essence
itself.} Just as the Ideas develop out of his infinite understanding by an
eternal birth, so the seed, the fundamental forms of all things, lies enveloped
in matter's womb, and only in time, and consequently only gradually, becomes
intuitable images of those patterns which they ever more nearly strive to
express.

The first human beings did not arise immediately either from the sea or from the
earth. Yet in both lies the mediate ground of both our matter and our form. Of
the earth we are formed. For before the solid separated itself from the fluid,
nothing living yet existed, least of all the most excellent creature among all
living things. Nature's forces are not different from the Godhead's own force,
except insofar as the latter considered in itself is infinite, while in the
former it reveals itself in a finite manner, and in its infinity only as an
unceasingly increasing magnitude. In this sense one may say that matter has
formed itself, advancing gradually from the less perfect to the most perfect
human form. One may likewise say with truth that God, after having brought forth
matter itself---or, by its properties, having made his own being as intuitable as
it could become---has out of it again brought forth all things. These turns of
speech, which most hold to be as different as heaven and earth, mean at bottom
the same, and indicate only two human manners of representation.

In a still more mediate way water is the origin of all things. The ancients,
who out of an easily excusable impatience leapt past the nearest and
intermediate causes in order
% printed p. 111
to soar up to the highest, misunderstood these propositions about our origin,
which are in themselves quite correct. Either they let human beings and gods
arise immediately out of chaos, which is the original fluid; or else they held
that by the united forces of the sun, the air, the water and the earth, the most
perfect animals and human beings themselves were begotten of the last, just as
mold and maggots are of putrefied bodies. The first opinion rested on a mythical
tradition and a poetic presentation of it. The poet's imagination resembles the
chaos they sing of, and in the same way contains the seed of reason's most
glorious works. It comprehends the whole at a glance, but does not notice that
what in this flash of genius---as in the divine understanding---exists all at
once, can become actual for the senses only gradually. The second hypothesis is
indeed a fruit of experience and reflection, but an unripe one. For although the
newest experiences seem in this to agree with the most ancient---that some among
the lowest plant and animal species can come forth in this way without the usual
propagation from their like, indeed even wholly new ones that did not exist
before---we have nevertheless not the least ground, from nature's regular course,
to suppose the same of the larger and more developed organisms in the two
kingdoms of living things. Such giant strides, as a sudden transition from the
lowest to the highest degree of development, are no less at variance with
philosophy's fundamental principles than with experience and analogy. For from
the concept of time and finitude the opposite manifestly follows. How was any
time possible without order and a continuous advance in which there can be no
leap? How can one think the infinite being expressed in finite forms except by
an eternal approximation to the former; or conceive time as eternity's image,
unless it is unbroken and uniform throughout?

From materials that seem raw, because one becomes aware of no form in them
there occurs, then, an imperceptible transition to more determinate shapes,
whose regular sequence or
% printed p. 112
series one can for that reason only indicate in general, as I have also done
above. To separate them by sharply marked boundaries will therefore not do,
because nature's course, in those deviations or degenerations by which the
former genera and species are transformed into wholly new ones, is not merely
creeping, but so slow withal that neither our own experience nor history's far
too small compass puts us in a position to follow its almost imperceptible
traces. After the lapse of millennia, the least ambiguous remains of animals and
human beings---which may be regarded as ancestors of those now existing, though
rather different---still give occasion for objections and doubts. I shall
therefore adduce none here. One will find them described in abundance in the
writings of naturalists, but most of them uncertain, all of them disputed. What
may be right or wrong with respect to these one cannot, as everyone sees, decide
\emph{a priori}; but the matter itself, which one would prove from them, is
nevertheless not doubtful.

An immediate creation, whether of actual animals and plants or of their nearest
germs, is empty words. If, on the other hand, the question is of a mediate or
natural production, then no other is possible than the one adduced. The human
race has consequently not existed from the beginning in its present form and
shape, although every history, the sacred as well as the secular, must naturally
proceed from this point. The specific criteria by which natural history now
distinguishes this being from others, it has not always had, but has taken on
gradually. I venture boldly to advance this opinion, not as a mere hypothesis,
but as an undeniable consequence of an order of nature that is not accidental
but necessary, whose series one cannot break off without absurdity and
contradiction. It would, on the other hand, be temerity if I were to maintain
that one can with the same certainty depict the human shape in every preceding
condition, or determine to which species of creatures now known we formerly
belonged, or from which we have perhaps sprung. I am convinced
% printed p. 113
that the human being constitutes a species or genus for itself, and by no means
believe that the ape---as the Tibetans, according to \emph{Pallas}'s report,
maintain---or any other animal, however high they may rise through all eternity,
will ever become human, or just the same as, even if more than, what we now are.
But I am indeed of the opinion that our first ancestors, whose memory no history
and no archive can have preserved, ran through a course of life in which neither
could the higher gifts of mind come to that maturity which they reach in the
beautiful and fully developed body we now bear, nor was mere sensibility provided
with all the requisite instruments for so many kinds of sensations, so fine and
so vivid. Nature's manifoldness, the inexhaustible wealth of Ideas in the
Godhead's bosom, their destination to an eternal and ever-progressing actuality,
does not permit us the thought that all beings are nothing but the same thing at
different degrees of the height they must ascend. What, then, does the
inequality among creatures of one species mean, which despite this individual
difference can each in its own way be equally perfect? Why should such
manifoldness not have place among species as well?

Have we perhaps reason to be ashamed of such an origin? This question can be
answered quite briefly. We need not go so far back in time to find stronger
grounds for our humiliation. How many years is it since each of us stood far
below the life of the animals, indeed of the meanest plant? Is it not more
important to think of what we now are, and shall one day quite certainly become,
than of what we were---as if the proud \emph{cedar} would look down with scorn
upon the young shoot that springs up at its roots, or we upon the fetus whose
heart has scarcely begun to beat in the mother's womb. Everywhere all that lives
has the same root: divine is our descent, eternal and infinite our destiny.

\begin{center}\rule{0.22\linewidth}{0.4pt}\end{center}

\bigskip
% printed p. 114
% [§ 2 introduces F. C. Sibbern; § 3 presumably Howitz. Not translated here.]

\end{document}
